\section{Product Measures} 

\subsection{Complete Measures}

  \begin{definition}[Complete Measure Spaces]
    A measure space $(X, \mathcal{M}, \mu)$ is \textbf{complete} if $\mathcal{M}$ contains all subsets of sets of measure $0$. 
  \end{definition}

  \begin{example}
    Listed. 
    \begin{enumerate}
      \item $(\mathbb{R}, \mathcal{L}, m)$ is complete. 
      \item $(\mathbb{R}, \mathcal{B}, m)$ is complete since we shows that the Cantor set (a Borel set of Lebesgue measure $0$), contains a subset that is not Borel. 
    \end{enumerate}
  \end{example}

  \begin{theorem}[Every Measure Space can be Completed]
    Let $(X, \mathcal{M}, \mu)$ be a measure space. Define $\mathcal{M}_0$ to be the collection of subsets $E$ of $X$ of the form $E = A \cup B$ where 
    \begin{enumerate}
      \item $B \in \mathcal{M}$, 
      \item $A \subset C$ for some $C \in \mathcal{M}$ for which $\mu(C) = 0$.\footnote{So we are basically splitting $E$ into a measure $0$ part $C$ and everything else $B$.}
    \end{enumerate}
    For such a set $E$, define $\mu_0 (E) = \mu(B)$. Then, $\mathcal{M}_0$ is a $\sigma$-algebra that contains $\mathcal{M}$, $\mu_0$ is a measure that extends $\mu$, and $(X, \mathcal{M}_0, \mu_0)$ is a complete measure space. 
  \end{theorem}
  \begin{proof}
    
  \end{proof}

\subsection{Premeasures} 

  Note that given a set function $\mu$ over $\mathcal{S}$, its Carathéodory extension $\overline{\mu}$ need not agree with $\mu$ for sets in $\mathcal{S}$. We would like to find adequate assumptions such that $\overline{\mu}$ is an extension of $\mu$. 

  \begin{definition}[Premeasure]
    Let $\mathcal{S}$ be a collection of subsets of $X$ and $\mu: \mathcal{S} \to [0, +\infty]$ a set function. Then, $\mu$ is called a \textbf{premeasure} if  
    \begin{enumerate}
      \item $\mu$ is finitely additive, 
      \item $\mu$ is countably monotone, 
      \item if $\emptyset \in \mathcal{S}$, then $\mu(\emptyset) = 0$. 
    \end{enumerate}
  \end{definition}

  \begin{theorem}[Premeasure Condition for Carathéodory Measure]
    Let $\mathcal{S} \subset 2^X$ and $\mu: \mathcal{S} \to [0, +\infty]$ a set function. In order for the Carathéodory measure induced by $\mu$ be an extension of $\mu$, it is necessary that $\mu$ is a premeasure.  
  \end{theorem}
  \begin{proof}
    
  \end{proof}

  So being a premeasure is a necessary but not sufficient condition, but if we impose on $\mathcal{S}$ a finer set-theoretic structure, this necessary condition is also sufficient. 

  \begin{definition}[Closure Under Relative Complements]
    A collection $\mathcal{S} \subset 2^X$ is said to be closed w.r.t. the formation of relative complements provided 
    \begin{equation}
      A, B \in \mathcal{S} \implies A \setminus B \in \mathcal{S}
    \end{equation}
  \end{definition}

  \begin{theorem}[Premeasure over Set Closured Under Relative Complements Induces Carathéodory Extension]
    Let $\mu: \mathcal{S} \to [0, +\infty]$ be a premeasure on $\mathcal{S}$ that is closed w.r.t. the formation of relative complements. Then, the Carathéodory measure $\overline{\mu}: \mathcal{M} \to [0, +\infty]$ induced by $\mu$ is an extension of $\mu$, called the \textbf{Carathéodory extension} of $\mu$. 
  \end{theorem}
  \begin{proof}
    
  \end{proof}

  However, a number of natural premeasures such as the premeasure length defined on the collection of bounded intervals of reals numbers, are defined on collections of sets that are not closed w.r.t. relative complements. So we consider alternate conditions for extending measures. 

  \begin{definition}[Semiring]
    A nonempty collection $\mathcal{S}$ of subsets of a set $X$ is a \textbf{semiring} if 
    \begin{enumerate}
      \item \textit{Closure under finite intersections}. 
        \begin{equation}
          A, B \in \mathcal{S} \implies A \cap B \in \mathcal{S}
        \end{equation}
      \item \textit{Disjoint decomposition of relative complements}. 
        \begin{equation}
          A, B \in \mathcal{S} \implies A \setminus B = \bigsqcup_{k=1}^n C_k
        \end{equation}
        for some collection $C_k \in \mathcal{S}$. 
    \end{enumerate}
  \end{definition}

  \begin{theorem}[Carathéodory-Hahn Theorem]
    Let $\mu: \mathcal{S} \to [0, +\infty]$ be a premeasures on a semiring $\mathcal{S}$ of subsets of $X$. 
    \begin{enumerate}
      \item Then, the Carathéodory measure $\overline{\mu}$ induced by $\mu$ is an extension of $\mu$. 
      \item Furthermore, if $\mu$ is $\sigma$-finite, then so is $\overline{\mu}$, and $\overline{\mu}$ is the unique measure on the $\sigma$-algebra of $\mu^\ast$-measurable sets that extends $\mu$. 
    \end{enumerate}
  \end{theorem}

\subsection{Product Measure} 

  Before, we saw how we can construct measures using Carathéodory construction. We will consider how to create product measures, which will be an extension of a set functions using the Carathéodory-Hahn theorem. 

  \begin{definition}[Measurable Rectangle]
    Let $(X, \mathcal{A}, \mu)$, $(Y, \mathcal{B}, \nu)$ be two measure spaces. Consider the product space $X \times Y$. If $A \in \mathcal{A}, B \in \mathcal{B}$, then the set $A \times B$ is called a \textbf{measurable rectangle}. 
  \end{definition}

  \begin{lemma} 
    Let $\{A_k \times B_k\}_{k=1}^\infty$ be a countable disjoint collection of measurable rectangles whose union is also a measurable rectangle $A \times B$. Then, 
    \begin{equation}
      \mu(A) \times \nu(B) = \sum_{k=1}^\infty \mu(A_k) \times \nu(B_k)
    \end{equation}
  \end{lemma}
  \begin{proof}
    
  \end{proof}

  We want to set up the conditions to invoke the Carathéodory-Hahn theorem. This naturally leads to the following. 

  \begin{theorem}
    Let $\mathcal{R}$ be the collection of measurable rectangles in $X \times Y$ and for a measurable rectangle $A \times B$, define 
    \begin{equation}
      \lambda(A \times B) = \mu(A) \cdot \nu(B)
    \end{equation}
    Then, $\mathcal{R}$ is a semiring and $\lambda: \mathcal{R} \to [0, +\infty]$ is a premeasure. 
  \end{theorem}
  \begin{proof}
    
  \end{proof}

  \begin{definition}[Product Measure]
    Let $(X, \mathcal{A}, \mu)$, $(Y, \mathcal{B}, \nu)$ be two measure spaces, $\mathcal{R}$ the collection of measurable rectangles contained in $X \times Y$, and $\lambda$ the premeasure defined on $\mathcal{R}$ by 
    \begin{equation}
      \lambda(A \times B) = \mu(A) \cdot \nu(B)
    \end{equation}
    for all $A \times B \in \mathcal{R}$. Then, the \textbf{product measure} $\lambda = \mu \times \nu$ is the Carathéodory extension of $\lambda: \mathcal{R} \to [0, +\infty]$ defined on the $\sigma$-algebra of $(\mu \times \nu)^\ast$-measurable subsets of $X \times Y$. 
  \end{definition}

\subsection{Fubini's and Tonelli's Theorem}

  \begin{theorem}[Fubini's Theorem]
    Let $(X, \mathcal{A}, \mu)$, $(Y, \mathcal{B}, \nu)$ be two measure spaces and $\nu$ be complete. Let $f$ be integrable over $X \times Y$ w.r.t. to the product measure $\mu \times \nu$. Then for almost all $x \in X$, the $x$-section of $f$, $f(x, \cdot)$, is integrable over $Y$ with respect to $\nu$, and 
    \begin{equation}
      \int_{X \times Y} f \, d(\mu \times \nu) = \int_X \bigg[ \int_Y f(x, y) \,d\nu(y) \bigg] \, d\mu(x)
    \end{equation}
  \end{theorem}

  \begin{theorem}[Tonelli's Theorem]
    Let $(X, \mathcal{A}, \mu)$, $(Y, \mathcal{B}, \nu)$ be two measure spaces and $\nu$ be complete.Let $f$ be a nonnegative $(\mu \times \nu)$-measurable function on $X \times Y$. Then, 
    \begin{enumerate}
      \item For almost all $x \in X$, the $x$-section of $f$, $f(x, \cdot)$, is $\nu$-measurable. 
      \item The function defined almost everywhere on $x$ by 
        \begin{equation}
          x \mapsto \int_Y f(x, y) \,d\nu(y)
        \end{equation}
        is $\mu$-measurable. 
      \item Finally, 
        \begin{equation}
          \int_{X \times Y} f \,d(\mu \times \nu) = \int_X \bigg[ \int_Y f(x, y) \,d\nu(y) \bigg] \,d \mu(x)
        \end{equation}
    \end{enumerate}
  \end{theorem}

\subsection{Exercises} 

  \begin{exercise}[Math 631 Fall 2025, Final Exam Exercise 4]
    Let $A=[-1,1]\times[-1,1]$. Let $f(x,y)=\frac{xy}{(x^{2}+y^{2})^{2}}$ (when $x=0$ or $y=0$ we set $f(x,y)=0$). Prove that the iterated integrals exist and are equal:
    \begin{equation}
      \int_{-1}^{1}(\int_{-1}^{1}f(x,y)dx)dy=\int_{-1}^{1}(\int_{-1}^{1}f(x,y)dy)dx,
    \end{equation}
    but the double integral $\int_{A}f(x,y) dxdy$ does not exist.
  \end{exercise}
  \begin{solution}
    Note that $f$ is odd with respect to both $x$ and $y$. Therefore
    \begin{equation}
      \int_{-1}^{1}f(x,y)dx = \int_{-1}^{1}f(x,y)dy=0
    \end{equation}
    for all $y$ and $x$ respectively. Hence both iterated integrals are zero. Now passing to polar coordinates, we have that $|f(r,\theta)|=\frac{|\sin 2\theta|}{2r^{2}}$. Therefore, for every $\epsilon>0$,
    \begin{equation}
      \int_{A}|f(x,y)|dxdy\ge\int_{\epsilon}^{1}\int_{0}^{\pi/2}\frac{\sin 2\theta}{r^{2}}rd\theta dr\ge c\int_{\epsilon}^{1}r^{-1}dr.
    \end{equation}
    Since $\int_{\epsilon}^{1}r^{-1}dr$ diverges as $\epsilon\rightarrow0$, $f$ is not integrable over $A$.
  \end{solution}



